\section{相关工作}

\subsection{多智能体强化学习}

多智能体强化学习已广泛用于合作控制、资源分配和对抗决策等任务。PPO 通过截断策略优化目标提高了 on-policy 策略梯度训练的稳定性~\cite{schulman2017ppo}。在多智能体场景中，MADDPG、COMA、VDN 和 QMIX 分别从集中式 critic、反事实信用分配和值函数分解等角度缓解非平稳训练和协同信用分配问题~\cite{lowe2017maddpg,foerster2018coma,sunehag2018vdn,rashid2018qmix}。集中训练、分散执行框架通过集中式 critic 或集中式价值分解利用全局信息，同时保留执行阶段的分散性。MAPPO 是合作式多智能体任务中的常用强基线。Yu 等人的研究表明，经过合适实现和超参数设置后，PPO 类多智能体方法可以在多个合作式基准任务中取得有竞争力的性能~\cite{yu2021mappo}。因此，本文将 MAPPO 作为严肃基线，而不是弱参考方法。

在无人机协同任务中，强化学习被用于轨迹规划、目标追逃、空战机动和集群协同。近期 UAV cooperative pursuit-evasion 研究也表明，强化学习可以学习复杂追逃策略~\cite{zhao2024uav_pursuit_evasion}。但在有限通信条件下，单纯依赖局部观测和集中式 critic 仍可能出现种子间方差较大、低通信半径下碰撞率较高等问题。

\subsection{图神经网络与多智能体协同}

图神经网络适合处理多智能体系统中的关系建模问题。GAT 通过 masked self-attention 让节点自适应聚合邻居信息~\cite{velickovic2017gat}。在多智能体强化学习中，MAGNet 等工作探索了图网络在深度多智能体强化学习中的应用~\cite{malysheva2020magnet}。近期综述进一步总结了 GNN 与 MARL 在通信和协同中的结合方式~\cite{liu2024gnn_marl,cuzin2026gnn_comm_survey}。

然而，普通图注意力通常主要基于节点嵌入计算注意力分数。对于无人机协同追逃任务，智能体间相对距离、方位、速度差和通信可达性是影响协同决策的关键变量。如果这些物理关系只隐含在节点状态中，策略需要自行推断边关系，学习难度较高。本文因此将相对边特征直接加入注意力得分，使图编码器显式感知空间关系和通信拓扑。

\subsection{有限通信无人机协同}

现实无人机系统中的通信受到距离、带宽、链路质量和遮挡等因素影响。IC3Net 等工作关注智能体何时通信以及如何在多智能体任务中学习通信门控~\cite{singh2018ic3net}。UAV 通信网络和移动边缘计算任务中，已有研究结合图注意力和多智能体强化学习进行轨迹设计、资源分配和协同优化~\cite{feng2024gat_uav_comm,kim2024uav_mec_madrl}。

与上述工作不同，本文关注的是协同追逃任务中的有限通信鲁棒性。本文显式设置通信半径，并在半径 4、6、8、10 下评估策略性能。实验不仅关注成功率，也关注碰撞率和种子间方差，以反映策略在安全性和稳定性方面的实际可用性。

\subsection{本文定位}

综上，现有工作分别从 MAPPO 类稳定训练、图神经网络关系建模和学习通信机制等角度推进了多智能体协同决策研究，但三者在有限通信无人机追逃任务中仍存在结合空间：第一，仅使用 MAPPO 难以显式表达智能体间的空间关系；第二，普通 GAT 能聚合邻居节点，但对距离、方位、速度差和通信可达性等边语义利用不足；第三，固定通信拓扑训练难以保证策略在通信半径变化时稳定。本文围绕这一缺口，将边特征增强角色图编码器与 MAPPO 训练框架结合，并通过分阶段随机通信半径微调提升跨半径鲁棒性。
